\documentclass[10pt,a4paper]{article}
\usepackage[margin=0.75in]{geometry}
\usepackage[T1]{fontenc}
\usepackage[utf8]{inputenc}
\usepackage{enumitem}
\usepackage{titlesec}
\usepackage{hyperref}
\hypersetup{colorlinks=true,urlcolor=blue!60!black,pdftitle={Pavan Raviteja Upadhyayula CV}}

\pagestyle{empty}
\setlist[itemize]{leftmargin=*,itemsep=1pt,topsep=2pt,parsep=0pt}
\setlength{\parindent}{0pt}
\titleformat{\section}{\large\bfseries}{}{0em}{}[\titlerule]
\titlespacing{\section}{0pt}{10pt}{4pt}

% \entry{Role}{Dates}{Place}
\newcommand{\entry}[3]{%
  \textbf{#1}\hfill{#2}\\\textit{#3}\vspace{2pt}}

\begin{document}

\begin{center}
{\LARGE\bfseries PAVAN RAVITEJA UPADHYAYULA}\\[4pt]
Saarbrücken, Germany $|$ \href{mailto:pavanteja0325@gmail.com}{pavanteja0325@gmail.com} $|$
\href{https://linkedin.com/in/pavan-upadhyayula}{linkedin.com/in/pavan-upadhyayula} $|$
\href{https://orcid.org/0009-0008-1235-4531}{orcid.org/0009-0008-1235-4531}
\end{center}

\section{Professional Summary}
Master's student in Cybersecurity at Universität des Saarlandes, researching user privacy perceptions and their mental models as part of my master's thesis. Background in empirical research methods and software security. Currently seeking PhD opportunities in the domain of Usable Security \& Privacy, with a primary focus on socio-technical issues in the Global South.

\section{Education}
\entry{MSc Cybersecurity}{04/2023 -- Present}{Universität des Saarlandes, Germany}
\begin{itemize}
\item Thesis submission: Nov 2026
\item Current grade: 2,3
\item Relevant coursework: Human-Computer Interaction, Research methods in HCI, Usable security
\item Santander-Stipendium | Stipendienprogramm UdS International - Sommersemester 2026
\end{itemize}
\vspace{4pt}
\entry{BTech Computer Science}{06/2018 -- 04/2022}{Koneru Lakshmaiah Educational Foundation (KL University), India}
\begin{itemize}
\item Final grade: 8.53/10 (German grade: 1,7)
\end{itemize}

\section{Skills}
\begin{itemize}
\item \textbf{Empirical Methods:} Survey design, mixed-methods analysis, non-parametric testing, thematic analysis
\item \textbf{Tools:} R, Python, ATLAS.ti, \LaTeX, Git, Go
\item \textbf{Focus Areas:} Usable privacy, Human-Computer Interaction, privacy by design, user mental models, privacy regulation
\end{itemize}

\section{Research Experience}
\entry{Master's Thesis Researcher}{11/2025 -- Present}{Internet Architecture (INET) group, MPI-INF, Saarbrücken}
\begin{itemize}
\item Conducting research on user privacy perceptions of their personal data and its processing in India, in light of its Digital Personal Data Protection Act (DPDPA)
\item Designed an exploratory survey (n=239) to understand the perspectives of the educated, English-speaking subset of the Indian population, considered the early adopters of such laws and their provisions
\item Analyzed Likert and open-text responses in a mixed-methods design: Friedman \& Mann-Whitney U tests, Spearman correlations, and ordinal/linear regression with Holm correction
\item Analyzed open-text responses analyzed via reflexive thematic analysis in ATLAS.ti
\item Identified the gap between users' privacy concerns, their varied personal data sensitivities, and the provisions of the law expecting uniformity without transparency
\item Found near-null demographic effects on users' personal data sensitivity and perceived control over its processing, indicating an underlying distrust towards current regulatory practices
\end{itemize}
\vspace{4pt}
\entry{Research Assistant}{05/2025 -- Present}{Fraunhofer IEM, Paderborn (Remote)}
\begin{itemize}
\item Identified a significant research gap in security metrics and automated risk assessment for software products through a comprehensive literature review
\item Evaluated the git repository of an IoT team at DB Systel to help boost their security posture
\item Extended the OSSF Scorecard tool (Go) with a new probe measuring pipeline reliability of the main branch
\item Applied the extended tool across 25+ repositories; reported concrete security-posture improvements
\end{itemize}

\section{Publication}
P. R. Upadhyayula, K. Amarendra and S. B. Kudupudi, ``Intrusion Detection of Imbalanced Network Traffic,'' \textit{2022 7th International Conference on Communication and Electronics Systems (ICCES)}, Coimbatore, India, 2022, pp. 840--844, doi: \href{https://doi.org/10.1109/ICCES54183.2022.9835996}{10.1109/ICCES54183.2022.9835996}.

\section{Teaching \& Outreach}
\entry{Student Peer Mentor}{04/2019 -- 05/2020}{KL University, India}
\begin{itemize}
\item Mentored 300+ peers in object-oriented programming, network security, and software engineering across 10+ interactive, hands-on learning sessions
\end{itemize}
\vspace{4pt}
\entry{Public Relations Chair}{03/2019 -- 11/2021}{White Hat Hackers Club, KL University}
\begin{itemize}
\item Led cybersecurity awareness initiatives, organizing network security workshops for 100+ students and running 4 CTF competitions
\end{itemize}

\section{Projects}
\textbf{Intrusion Detection of Imbalanced Network Traffic} (Bachelor's Thesis)
\begin{itemize}
\item Designed and implemented an intrusion detection system leveraging the Modern Random Forest algorithm
\item Conducted metadata analysis across multiple KDD datasets, achieving a false alarm rate of $\sim$5\%
\end{itemize}

%\section{Certifications}
%\begin{itemize}
%\item GIAC Security Essentials (GSEC), SANS Institute --- workshop completed, certification in progress
%\end{itemize}

\section{Languages}
\begin{itemize}
\item English: Fluent (C1) \quad $|$ \quad German: Basic (A2)
\end{itemize}

\end{document}